% Generated from paper/full_draft. Edit the Markdown source or migration script.

Retrieval-augmented systems must select enough evidence to support an answer while limiting noise and language-model context cost. We formulate this as a two-stage route-control and budget-calibration problem. IntentRoute combines dense retrieval, BM25, geometry-defined cluster-local routes, and trust-weighted LinUCB feedback. Route confidence controls routing and fallback, while an independently calibrated policy sets the final context budget. Dense retrieval remains a recall floor. A bounded piecewise relevance-manifold hypothesis motivates local route construction; geometry is evaluated as a diagnostic signal rather than a standalone retrieval theory.

On LoTTE technology/search from 100k to 638k chunks, eligible frozen policies reduce evidence-context tokens by 6-18\% while avoiding the larger $\mathrm{Hit@10}$ losses of dense-only adaptive truncation; the original 400k split remains calibration-ineligible. A normalized five-fold 400k follow-up yields 14.50\% mean saving with no mean Hit change, although strict seed-level non-inferiority remains unestablished. Matched BGE-base and E5-base tests retain near-dense $\mathrm{Hit@10}$ with about 12\% token reduction. Route controls show that geometry and feedback improve route-level quality, but do not establish safe per-query compression without calibration and rescue. In a frozen 300-query downstream evaluation, matched variants reduce context by 6-12\%. DeepSeek, GLM-5.2, MiniMax-M3, and majority-vote comparisons show no statistically detectable correctness difference, while faithfulness remains method-dependent and degrades for the tested BGE policy. Prompt compression and reranking remain complementary. The supported contribution is a geometry-guided, feedback-adaptive route controller with separate budget calibration, not universal superiority over dense retrieval or proof that geometry determines relevance.
